\documentclass[10pt,a4paper]{amsart}
\usepackage[margin=19mm]{geometry}
\usepackage{amsmath,amssymb,mathtools}
\usepackage[scr=rsfs,cal=euler]{mathalfa}
\usepackage{xfrac}
\usepackage[unicode,hidelinks,bookmarks=false]{hyperref}
\newcommand{\ie}{i.e.\ }
\newcommand{\cf}{cf.\ }
\newcommand{\resp}{respectively\ }
\newcommand{\Subsection}{\subsection}
\providecommand{\nobreakdash}{\nobreak\hbox{-}}
\setlength{\emergencystretch}{2em}
\multlinegap0pt
\usepackage[shortlabels]{enumitem}
\usepackage{amssymb}
\usepackage{dsfont}
\usepackage{mathtools}
\usepackage[normalem]{ulem}
\usepackage{xcolor}
\usepackage{tikz}
\newtheorem{proposition}{Proposition}
    \def\R{\mathbb{R}}
    \renewcommand\geq{\geqslant}
    \newcommand{\rank}{\mathrm{rank}}
\title{Lifting curved Kakeya sets to linear Kakeya sets}
\author{Arian Nadjimzadah}
\date{Source: arXiv:2609.06299v1 (CC BY 4.0)}
\begin{document}
\maketitle

\noindent\emph{Source and licence.} Lifting curved Kakeya sets to linear Kakeya sets. Version arXiv:2609.06299v1 (CC BY 4.0), distributed by arXiv under the Creative Commons Attribution 4.0 International licence, http://creativecommons.org/licenses/by/4.0/, as recorded in the arXiv licence metadata for this exact version and shown on the abstract page https://arxiv.org/abs/2609.06299v1. Full licence notice: \texttt{licenses/arxiv-2609.06299-CC-BY-4.0.txt}.\par\smallskip

\noindent\textit{Editorial scope.} Counted unit: the article's proposition examples (source0.tex lines 476--486). The article's complete original proof of it is delivered with it (source0.tex lines 591--625, one \texttt{proof} environment of 35 lines, which the article places away from the statement). No mathematics is written, repaired or re-derived: the statement and the proof are the article's own lines, in the article's own notation, and no retained line is substituted or re-flowed.  Retained uncounted as the source-local context the proof needs: lines 345--347; lines 377--380; lines 438--440.  Nothing was omitted from any retained span, so the package carries no omission record.  No retained line contains a citation, so the wrapper carries no bibliography and no bibliography entry.  No retained line contains a figure environment, an image-inclusion command or drawing code, so no image is delivered and the package declares no asset dependency.  Editorial formatting only, in addition: an AMS \texttt{amsart} preamble that restates the article's own declarations and macros used by the retained lines, namely the environments \texttt{proposition} on separate counters, together with the standard packages those lines need.  The retained environments print in this document as Proposition 1 in order of appearance, whereas the article prints the same items with the numbering of its own full document.  The complete native source of the article as distributed in the arXiv e-print for this exact version is bundled unchanged as \texttt{sources/209519/source0.tex}.
    \begin{align}\label{eq: quadratic def}
        X_B(w,t;y) = (w + (tI + t^2 B)y).
    \end{align}

    \begin{align}
        &X^{(n)}_{\tan} : [-0.1,0.1]^{n-1} \times [0.9,1.1] \times [-0.1,0.1]^{n-1} \to \R^{n-1}, \\
        &X_{\tan}(w,t;y) = (w'-t^2 y', \tan^{-1}(w_{n-1}/t)-ty_{n-1}), \label{eq: tan example def}
    \end{align}

            \begin{align}\label{eq: projection property}
                \pi(L_{y,w,\rho}(t)) = \gamma_{y,w}(t).
            \end{align}

    \begin{proposition}[Examples of lifts of families of curves] \label{prop: lifting examples}\hspace{0pt} 
        \begin{enumerate}[label=(\roman*)]
            \item (Quadratic) Fix $n \geq 3$. The defining function $X_B$ from \eqref{eq: quadratic def} has a $(\rank B, \rank B - \rank B^2)$-lift.
            \item (Polynomial) Fix $n\geq 3$, $D \geq 1$, compact balls $W,Y \subset \R^{n-1}$, and a compact interval $I \subset \R$. Consider the defining function $X_D : W \times I \times Y$ given by 
            \begin{align}
                X_D(w,t;y) := A_0(w,y) + tA_1(w,y) + \cdots + t^{D-1} A_{D-1}(w,y) + t^D B(y),
            \end{align}
            where the $A_i$ are smooth functions, $\det DB \neq 0$, and $X_D$ is H\"ormander-type. Then $X_D$ has a $((D-1)(n-1), 0)$-lift.
            \item ($\tan$-example) The defining function $X_{\tan}$ from \eqref{eq: tan example def} has an $(n,0)$-lift. 
        \end{enumerate}
    \end{proposition}

    \begin{proof}[Proof of Proposition \ref{prop: lifting examples}(i)]
        Let $B$ be an $(n-1)\times (n-1)$ matrix, let $(s,\nu) = (\rank B, \rank B - \rank B^2)$, and define $q = \rank B$. Consider the rank factorization 
        \begin{align}
            B = UV, \quad U : \R^q \to \R^{n-1} \text{ injective}, \quad V : \R^{n-1} \to \R^q\text{ surjective.}
         \end{align}
         Consider the coordinates $(x,z,t) \in \R^{n-1} \times \R^q \times \R = \R^{N}$ and define 
         \begin{align}
             \pi(x,z,t) = (x + tUz,t) \in \R^{n-1} \times \R.
         \end{align}
         This is a submersion because $D\pi$ has rank $n$. Let $R \subset \R^q$ be a compact ball. For $\rho \in \R$, set 
         \begin{align}
             a(y,w,\rho) = (w,\rho), \quad \Theta(y,\rho) = (y- U\rho,Vy).
         \end{align}
         We may compute 
         \begin{align}
             \pi((a,0) + t(\Theta,1)) &= (w + t(y-U\rho) + tU(\rho + tVy),t) \\
             &= (w + (tI +t^2B)y,t),
         \end{align}
         so \eqref{eq: projection property} holds.

         Now we show that $\rank D\Theta = N - 1 - \nu = (n-1) + \rank B^2$. The kernel of $D \Theta$ consists of pairs $(h,k)$ satisfying
         \begin{align}
            h - Uk = 0, \quad Vh=0,
         \end{align}
         so $\ker D\Theta$ is naturally identified with $\ker(VU)$. Thus $\dim\ker (VU) = \dim \ker(D\Theta)$.
         Since $U$ is injective and $V$ is surjective, 
         \begin{align}
             \rank(VU) = \rank(UVUV) = \rank(B^2).
         \end{align}
         Then by applying rank-nullity to $VU$ and to $D\Theta$, we obtain 
         \begin{align}
             \rank D\Theta &= (n-1) + \rank(B^2) \\
             &= N-1-\nu.
         \end{align}
    \end{proof}

\end{document}
